\chapter*{序言}
\addcontentsline{toc}{chapter}{序言}
\markboth{序言}{序言}

\section*{封面：周转圆的宇宙}

封面上，三个同心圆静静地嵌套着。最外圈的圆环上，六十个点均匀分布，如同钟表的刻度，暗示着六十甲子的循环；向内一层，十二个红点标记着地支的方位；再向内，八个绿点对应八卦的排列；最内圈，五个橙点暗示着五行的运转。圆心处，一个红点延伸出细线，连接着第一个小圆；小圆上的点又延伸出更小的圆；更小的圆上还有更小的圆。这是一幅周转圆的图景——外圈每前进一步，内圈旋转一整圈。古希腊人用这样的嵌套圆描述行星的运动，而本书将用它描述\textbf{在适当条件下}具有周期结构的事物。

封面的底部有两行公式。第一行告诉我们，一个周期为 $T$ 的系统，其相位点可以用三元组 $(A_k, \omega_k, \phi_k)$ 来标记——振幅、频率、相位。第二行则表达了周期卷积的初步思想：两个函数之间的深层相似性，不仅取决于各自频率是否相同，还取决于全部频率对之间的交互。\textbf{需要强调}：这一表述目前是一幅\textbf{启发性图景}，而非数学定理；其严格成立所需的条件见\textbf{引言}中周期卷积定义之后所列的收敛性、对称性与有界性要求。

这幅图景是本书全部内容的视觉概括：周期系统可以由周转圆族描述，周转圆由振幅、频率、相位刻画，周期之间的关系由周期卷积描述。\textbf{需要再次强调}：这一概括是一幅\textbf{研究纲领式的图景}，它在本书中被逐步转化为若干精确的、有条件的数学命题，而非一蹴而就的定理。

\section*{周期：人类理解世界的一种早期语言}

人类对周期的觉察，可以追溯到对天穹的观测。昼夜交替、月相盈亏、星辰回归——这些天文节律以可预见的规律重复出现，在多个早期文明中构成了时间意识的起点。在美索不达米亚，苏美尔人将金星的出没编成星表；在尼罗河谷，埃及人借天狼星的偕日升预报洪水；在黄河两岸，先民以圭表测量日影的消长，定节气，授农时。周期，是人类理性试图刻画时间的一种早期工具。

天穹之上的周期不仅催生了历法，也在若干文明中促进了数学的发展。当先民发现日月运行并非孤立现象，而是可以相互参照、彼此预测时，一种认识便逐渐成形：天体的运行在一定精度内可以由节律性的数学模型把握。这一认识，是精密科学传统在若干文明中独立出现的因素之一——但并非唯一因素，也并非所有文明都以相同方式展开这一认识。

\section*{周期的数学化}

周期一旦被察觉，便逐渐被数学化。巴比伦人将日月五星的运行周期化为六十进制的数表；玛雅人用多重历法嵌套追踪金星与朔望的交错；古希腊人将天体运动分解为本轮与均轮的复合。在若干早期文明中，天文学是最早达到高精度的科学分支之一——因为天体的周期性运动提供了相对稳定的测量对象。

在中国，这一数学化进程以独特的方式展开。《淮南子·天文训》以“七十六岁，日月复以正月入营室五度，无余分，名曰一纪”确立了回归年与朔望月的一种高精度公倍近似；《皇极经世》以元会运世的嵌套体系——一元十二会，一会三十运，一运十二世，一世三十年——将天象周期与人事周期纳入同一个自洽的算术框架；二十四节气的划分则是太阳历的一种独立创造：将黄道一周等分为二十四个相位节点，每个节点对应太阳在黄道上的近似位置。\textbf{需要指出}：上述历史叙述属于传统历法推演与文化史解释，其具体数值与年代设定应视为\textbf{象征性设定或据传统说法}，不宜作为严格历史事实或数学推导的前提。

到了十七世纪，开普勒将第谷二十余年的观测数据化为三条定律。行星的椭圆轨道不再是任意设定的几何构造，而是在平方反比力场中角动量守恒与能量守恒共同约束下的必然结果。牛顿进一步将这三条定律统一为万有引力方程，周期性的数学描述从此获得了经典力学的根基。\textbf{需要指出}：开普勒定律与牛顿力学对周期运动的解释，是\textbf{在经典力学框架内}成立的；在广义相对论或量子力学框架中，相应的表述需要修正。

\section*{从不变性到守恒律，从守恒律到对称性}

周期性的一个核心特征是\textbf{不变性}。日出日落，今天和昨天相似；月缺月圆，这一轮和上一轮相似；春去秋来，这一年和上一年相似。先民首先觉察到的，是某些属性在变化中保持近似不变。太阳每天升起的方向大致相同，月亮每二十九天左右完成一次圆缺，四季每年按固定顺序更替。这些近似的不变性，是最早的周期知识。

当不变性被精确化为数学命题时，守恒律便登场了。一个量在某一类演化过程中保持不变，我们就说这个量在该类过程中守恒。开普勒发现行星在椭圆轨道上运动时，行星与太阳的连线在相等时间内扫过相等的面积——这是角动量守恒的一个特例。伽利略发现摆的周期在小振幅下与振幅近似无关——这是非线性系统在小振幅极限下等时性的表现，并非严格的能量守恒。牛顿将这些观察统一为运动定律：在平方反比力场中，角动量守恒与能量守恒共同约束轨道为闭合椭圆。

守恒律一旦被确立，对称性便自然浮现。一个系统如果具有守恒量，就意味着存在某种变换，在这种变换下系统的行为保持不变。时间平移不改变能量，空间平移不改变动量，旋转不改变角动量。诺特定理在二十世纪初为这一直觉提供了深刻的数学形式：\textbf{在作用量原理成立、且对称变换是连续光滑的前提下}，每一个连续对称性蕴含一个守恒量，即\textbf{对称 $\Rightarrow$ 守恒}；\textbf{反之不成立}——守恒律并不蕴含相应的对称性（见「周期基础」章 \cref{cor:conservation_symmetry_hierarchy}）。\textbf{诺特定理的条件不可省略}——对离散对称性、非拉格朗日系统或非光滑系统，上述蕴含一般不成立，需要分别处理。

需要强调逻辑顺序：\textbf{不是先有对称性再有守恒律，而是先有不变性的观察，再抽象出守恒律，再发现对称性}。在诺特定理的框架下，守恒律与对称性并非一一对应，而是\textbf{单向蕴含}：\textbf{对称 $\Rightarrow$ 守恒}，而\textbf{守恒 $\nRightarrow$ 对称}（\cref{cor:conservation_symmetry_hierarchy}）——两者共同指向的事实是：系统在变化中保持着某种恒定的结构。这一结构在相空间中的表现，\textbf{在一定条件下}就是闭合轨道或不变环面。\textbf{但守恒律本身并不保证轨道闭合}——一个具有能量守恒的系统，其轨道也可能是不闭合的（如双曲轨道）；闭合还需要额外的拓扑或动力学条件。

\section*{周期性的三种生成模式：一个研究纲领}

在本书的框架中，我们\textbf{将周期性的生成变换归结为三种基本操作}，并主张它们构成一个合理的研究纲领。这一主张的严格形式——即“任何保持某种内禀结构的连续变换都可分解为这三种群的复合”——是一个\textbf{研究纲领}，其成立范围需要逐章限定，目前不应视为已证定理。

\begin{enumerate}[leftmargin=2em, label={(\arabic*)}]
    \item \textbf{平移}。一辆匀速行驶的火车上，车轮每旋转一周，车厢地板上某个标记便回到相同位置。平移生成的是传统的循环周期，其\textbf{在拉格朗日框架下的}守恒量为动量。
    \item \textbf{缩放}。科赫雪花的边界在每一次缩放变换下都回归自身——放大三倍后，边界曲线与放大前结构相同。这是一种\textbf{自相似}结构：周长无限而面积有限，面积正是使缩放“闭合”的守恒量。
    \item \textbf{旋转}。地球绕太阳的公转，在忽略其他行星摄动的近似下，是一个旋转对称系统中的闭合轨道，其守恒量为角动量。
\end{enumerate}

梯度、旋度、散度三个算子在「矢量分析」章中各自建立严格定义、可微性前提与适用范围；它们各自的结构已在\textbf{引言}中列出，此处不再重复。

\textbf{需要特别指出}：本书\textbf{不}把这三个算子与平移、缩放、旋转三类基本变换并列为一组“对应关系”。三类变换作用在\textbf{空间}上，三个算子作用在\textbf{场}（标量场、向量场）上，二者所处的数学层次不同；把它们排成一张对应表既不必要，也容易掩盖各自的真实前提。凡涉及三者中任何一条的陈述，均在「矢量分析」章中以明确的可微性前提与定义域条件单独给出。

\section*{相位：关系本身的数学}

在周期数学中，相位\textbf{可以}被理解为主动的变换算子，而不仅仅是被动的坐标。周期 $T$ 的第 $k$ 个相位对应复数 $\mathcal{O}_k^{(T)} = e^{i2\pi k/T}$。乘以 $\mathcal{O}_0 = 1$ 得到自身——合同；乘以 $\mathcal{O}_{T/2} = -1$ 得到反向量——对径；乘以任意 $\mathcal{O}_k$ 定义旋转 $2\pi k/T$ 的关系。相位算子的乘法封闭性保证了关系体系的封闭性：关系之间的关系仍然是关系。这一封闭性是有限循环群 $C_T$ 的代数性质，不依赖于任何物理诠释。

这一视角为将《易经》中的“合化”“冲克”等传统关系纳入群表示论框架提供了一种\textbf{可能的}路径；其具体的现代数学编码（天干、地支、八卦与具有周期轨道的群的对应关系，以及其中的周转圆嵌套规则）见\textbf{引言}，本书在那里把这些对应表述为代数同构或群作用等价，而不是对古代文献的历史断言。

\section*{莱布尼茨的发现：周期内嵌于数本身}

在周期的数学化进程中，莱布尼茨做出了一个独特而深刻的观察。他在研究二进制算术时发现，自然数序列的二进制展开中，每一位上的数字呈现出严格的周期性——第 $0$ 位以 $2$ 为周期循环，第 $1$ 位以 $4$ 为周期循环，第 $k$ 位以 $2^{k+1}$ 为周期循环。

\textbf{历史考证提示}：莱布尼茨曾研究二进制算术并对其哲学意涵（特别是与《易经》的关联）深感兴趣；本书中“他将这一现象命名为 \textbf{Periodus numerorum}——自然数的周期”一句，属于\textbf{待考史料}，需要原始文献或二手研究的支持才能作为历史事实陈述。此处我们保留这一提法，但标注其认识论地位为“据称”或“史料待考”，不作历史定论。

这一观察的\textbf{数学内容}本身是清楚的：周期性并非仅存在于天穹之上的宏大运动之中，也内嵌于数系的二进制表示结构之中。每一位上的 $0$ 与 $1$ 以固定的节律交替。

莱布尼茨更进一步指出，二进制小数在复平面单位圆上的均匀分布——取 $n$ 位二进制小数可将圆周等分为 $2^n$ 个点——构成了对连续圆周的一种逼近。这一几何直观\textbf{可以}与《易经》卦象体系的生成结构相对应：太极（$2^0$）$\to$ 阴阳（$2^1$）$\to$ 四象（$2^2$）$\to$ 八卦（$2^3$）$\to$ 六十四卦（$2^6$）。\textbf{需要指出}：这种对应是本书给出的现代数学诠释，而非对《易经》原始文本的历史断言。

\section*{周转圆：描述周期现象的一种统一语言}

在具有守恒律的系统中，\textbf{在一定条件下}存在闭合轨道。一个具有角动量守恒的行星，其轨道在平方反比力场中闭合；一个具有能量守恒的单摆，其摆动在相空间中是一个闭合轨道。闭合轨道意味着系统在相空间中画出一个圈——周转圆。

周转圆不是天文学的专利。\textbf{许多}具有周期结构的事物，都可以用周转圆来描述。一束单色光的传播，可以用电磁场在空间中的正弦振荡描述；一个词语的声波振动，可以用空气压力的周期起伏近似描述；一个句子的韵律，可以用音节长短与重音的周期交替近似描述；一个电信号在导线中的振荡，可以用电荷在平衡位置附近的往复近似描述；一块芯片中的晶振，是一个基本固定的周转圆，为整个系统提供时间基准。

从数学上看，\textbf{在适当的函数空间和收敛条件下}，周期函数可表示为傅里叶级数：
\[
f(x) = \sum_{k \in \mathbb{Z}} c_k e^{i k \omega_0 x}.
\]
\textbf{需要区分两种情形}：(i) \textbf{有限谐波模型}：$f(x) = \sum_{k=1}^{N} A_k e^{i(\omega_k x + \phi_k)}$，适用于描述有限多个周转圆的叠加，是数值实现和工程建模中的常用形式；(ii) \textbf{一般傅里叶级数}：$f(x) = \sum_{k \in \mathbb{Z}} c_k e^{i k \omega_0 x}$，是对周期函数（在适当收敛性条件下）的完整展开，其系数为复数，包含振幅与相位信息。\textbf{不能说“周期函数 = 有限个周转圆的叠加”}——一般的周期函数需要无限多个傅里叶模式（例如方波需要无限多奇次谐波）。但当注意力集中在有限多个主导分量时，有限谐波模型是一个有用的近似。本书在涉及具体数值实现时通常采用有限谐波模型，而在涉及一般理论时使用傅里叶级数。

在适当的条件下，周期现象可以被理解为若干周转圆的叠加。外圈的大圆承载着整体的节奏，内圈的小圆贡献着细节的纹理，而相位决定了每个圆在某一时刻的具体位置。

\section*{周期卷积：跨周期交互的一种统一计算}

本书提出的一个核心研究纲领是\textbf{周期卷积}：把振幅（缩放）与相位（旋转）也纳入级数运算，从而保留全部频率对之间的交互；封面上第二行公式即其初步形态。它的两层构造（交互张量、合成映射）、由二者复合而成的周期卷积、以及严格成立所需的各项条件，统一在\textbf{引言}中给出，此处不重复。

\section*{逼近、置闰与不可公度周期的有界化}

在周期数学面对的众多技术困难中，有一类困难与众不同：它不来自计算、不来自拓扑，而来自\textbf{数论}。当连续周期 $T_f$ 与采样间隔 $\Delta x$ 不可公度，即理想周期比
\[
\rho = \frac{T_f}{\Delta x} \notin \mathbb{Q}
\]
时，任何有限窗口内的采样序列都\textbf{不可能}具有严格整数周期。这不是观测误差，不是仪器缺陷，而是不可公度性的内在代价。这种代价在频域中的表现就是\textbf{频谱泄漏}：窗口内周期数不是整数，信号能量必然从主瓣扩散到旁瓣。

\subsection*{置闰的数学定义}

本书提出的一个核心方法论是\textbf{置闰}——用一列整数 $\{N_m\}$ 逼近实数 $\rho$，使\textbf{累计采样数}
\[
S_M = \sum_{m=0}^{M-1} N_m
\]
以\textbf{有界误差}逼近 $M\rho$。置闰的核心构造是三种基本置闰序列：
\begin{enumerate}[leftmargin=2em, label={(\arabic*)}]
    \item \textbf{下置闰}：$N_m^{\mathrm{floor}} = \lfloor (m+1)\rho \rfloor - \lfloor m\rho \rfloor$，误差 $E_M^{\mathrm{floor}} \in (-1,0]$；
    \item \textbf{上置闰}：$N_m^{\mathrm{ceil}} = \lceil (m+1)\rho \rceil - \lceil m\rho \rceil$，误差 $E_M^{\mathrm{ceil}} \in [0,1)$；
    \item \textbf{平衡置闰}：$N_m^{\mathrm{round}} = \operatorname{round}((m+1)\rho) - \operatorname{round}(m\rho)$，误差 $|E_M^{\mathrm{round}}| \le 1/2$。
\end{enumerate}
三种序列的共同特征是：$N_m \in \{\lfloor\rho\rfloor, \lceil\rho\rceil\}$——每一步只在下整数周期与上整数周期之间选择——但累计误差被约束在 $[-1/2, 1/2]$、$(-1,0]$ 或 $[0,1)$ 的有界带内。这一有界性是置闰的全部价值所在。

\subsection*{逼近论基础：连分数的角色}

置闰序列的\textbf{最优}选择不是任意的，它由 $\rho$ 的连分数展开决定。设 $\rho$ 的连分数收敛子为 $\{p_j/q_j\}$。则：
\begin{enumerate}[leftmargin=2em, label={(\roman*)}]
    \item \textbf{最佳性}：$p_j/q_j$ 是所有分母不超过 $q_{j+1}-1$ 的有理数中最逼近 $\rho$ 者；
    \item \textbf{误差上界}：$|p_j - q_j\rho| < 1/q_{j+1}$；
    \item \textbf{交错逼近}：偶数项收敛子从下方、奇数项从上方交替逼近 $\rho$，形成对称夹逼。
\end{enumerate}
在置闰窗口 $T = p_j$、$M = q_j$ 上，累计误差趋于零。这使置闰从一个“误差控制技巧”提升为一个\textbf{有严格最优性判据的方法}——连分数收敛子就是最优置闰窗口。这一最优性来源于数论中一个经典事实：连分数收敛子是实数的最佳有理逼近；置闰把它从数论翻译为采样窗口选择。

\subsection*{对频谱泄漏的有界化}

置闰对频谱泄漏的作用可以精确表述。设 $F[k]$ 为置闰窗口上的 DFT，$c_m$ 为连续傅里叶系数。则
\[
\left| F[k] - c_{k/M} \right| \le C \cdot \frac{|E_M|}{T} \cdot \sum_{m\in\mathbb{Z}} |c_m|,
\]
其中 $C$ 为与 $T$ 无关的常数。这一不等式给出泄漏的\textbf{显式上界}——置闰把泄漏从“不可控”降为“有界”。当 $\rho \in \mathbb{Q}$ 时，置闰退化为常数列，泄漏严格为零；当 $\rho \notin \mathbb{Q}$ 时，置闰通过连分数选择窗口，使该上界在收敛窗口序列上趋于零。

\textbf{需要强调}：置闰\textbf{不消除}泄漏，而是把泄漏幅度从 $O(1)$ 降为 $O(|E_M|/T)$。这一区别是本质的，它反映了不可公度性的内在代价——任何有限窗口都无法使 $\rho$ 与整数精确对齐。置闰的价值在于：它给出了这一代价的\textbf{定量刻画}和\textbf{系统化的最小化方法}。

\subsection*{置闰与窗函数的互补}

置闰与窗函数是抑制频谱泄漏的两条独立路线，作用于泄漏的不同环节：
\begin{itemize}[leftmargin=2em]
    \item \textbf{置闰}作用于\textbf{窗口长度}：选择 $T = p_j$ 使窗口内近似包含整数个连续周期，从源头减少泄漏；
    \item \textbf{窗函数}作用于\textbf{窗口边界}：通过平滑窗（Hann、Hamming、Blackman 等）压旁瓣，代价是主瓣展宽。
\end{itemize}
两者可叠加使用：先用置闰选择最优窗口，再施加窗函数，可同时获得窄主瓣与低旁瓣。这是单一方法无法达到的组合效果。

\subsection*{置闰在本书中的位置}

置闰方法在本书中承担三项功能：
\begin{enumerate}[leftmargin=2em, label={(\arabic*)}]
    \item 在「离散采样逼近与置闰」章中，置闰被建立为公度性与非公度性之间的桥梁，给出置闰误差有界定理与连分数最佳置闰定理；
    \item 在「调和分析」章中，置闰被用于推导置闰窗口的 DFT 泄漏估计，给出泄漏幅度的显式上界；
    \item 在「周期范畴」章中，置闰被函子化——置闰函子 $\mathcal{R}$ 修复连续范畴到离散范畴的不可逆缺口，使连续极限与离散化函子的伴随对在非公度情形下保持到任意精度。
\end{enumerate}
此外，置闰与「时序分析」章中的窗函数方法、「天文系统」章中的历法置闰、「历法系统」章中的无中气置闰、「密码系统」章中的状态周期、「智能体」章中的多智能体同步、以及「数据分析」章中的多窗谱估计，构成一个贯穿全书的统一视角。

\textbf{需要再次强调}：置闰目前在全书中被建立为一项有严格定理支撑的方法论，但它\textbf{不解决}不可公度性的所有困难。多频信号、时变周期、未知周期比等情形下，置闰的作用受限，需要与其他方法（多窗谱估计、贝叶斯谱分析、联合连分数逼近）结合使用。本书将\textbf{明确标注}置闰方法的适用范围和局限性。

\section*{从学科分裂到统一视角}

在现有的学科体系中，不同领域的研究者都在处理周期现象，却使用着彼此割裂的语言。天文学家使用摄动理论与轨道共振描述天体的准周期运动；凝聚态物理学家使用空间群与布里渊区描述晶体的平移对称性；生物学家使用极限环与相位响应曲线描述昼夜节律的基因网络；分形几何学家使用自相似性与标度指数描述跨尺度的不变结构；经济学家使用谱分析与自回归模型探测商业周期的特征频率；计算机科学家使用循环冗余校验与快速算法处理信号的周期分量。

每一位研究者都在处理与周期性有关的问题，却使用着彼此割裂的术语体系。一个凝聚态物理学家说的“倒格矢”，一位信号处理工程师说的“奈奎斯特频率”，一位调和分析学家说的“庞特里亚金对偶”——它们\textbf{在有限循环群或离散平移群的调和分析中}指向同一数学事实：周期函数的频域表示。\textbf{需要指出}：这种“同一”是在适当的抽象层次上成立的，具体的物理意义与适用条件在不同学科中可能不同。

类似的割裂同样存在于积分变换领域。傅里叶变换、拉普拉斯变换、小波变换、梅林变换通常被作为彼此独立的工具分别教授。然而，从群表示论的角度审视，\textbf{若选择适当的群作用、表示和函数空间}，不同积分变换可以在群表示与调和分析的框架下获得统一描述：傅里叶变换与平移群相关，梅林变换与缩放群相关，小波变换与仿射群相关。\textbf{这些陈述需要逐例验证}，其精确形式取决于所选的函数空间、群作用与测度。

在采样理论中，这种割裂同样存在。“采样定理”在信号处理教科书、“庞特里亚金对偶”在调和分析教科书、“置闰”在历法天文学教科书——三者分别教授，却指向同一个数学结构的三个侧面：不可公度情形下如何用离散逼近连续。本书的置闰理论正是把三者纳入同一框架的尝试。

\section*{学科壁垒的代价}

上述术语与方法的割裂至少带来三重代价：

\begin{enumerate}[leftmargin=2em, label={(\arabic*)}]
    \item \textbf{重复发现}。某一学科已知的数学定理，在另一学科被独立重新发现，造成学术资源的重复投入。
    \item \textbf{深层类比被遮蔽}。开普勒椭圆与科赫雪花之间、晶体点阵与基因振荡之间、傅里叶变换与梅林变换之间——这些\textbf{在适当抽象层次上}成立的同一数学结构，因学科语言的差异而难以被识别。
    \item \textbf{缺乏诊断工具}。面对新的复杂系统，研究者可能只能依赖现象层面的描述，而难以进行结构层面的诊断。
\end{enumerate}

\section*{统一框架的可能性}

尽管学科表面各异，但在每一个涉及周期结构的案例中，\textbf{通常}可以分解出三个共同的数学构件：

\begin{itemize}[leftmargin=2em]
    \item 一个\textbf{对称群} $G$，作用于系统的状态空间 $M$ 上；
    \item 一个\textbf{守恒量} $Q$，在 $G$ 的作用下保持不变；
    \item 一个\textbf{闭合轨道} $\mathcal{O} \subset M$，由守恒律与动力学共同约束生成。
\end{itemize}

这意味着，一个以群作用为语言、以守恒量为判据、以闭合轨道为几何图像、以周转圆级数为分解工具的\textbf{统一周期数学框架}，在理论上是可能的。本书正是这一可能性的一次探索。\textbf{其成功与否，取决于它能否在后续各章中产生非平凡的新数学结果}，而不仅仅是对已有知识的重新命名。

本书将在后续章节中逐步展开置闰方法的完整理论：它在「离散采样逼近与置闰」章中作为公度化工具，在「调和分析」章中作为泄漏有界化工具，在「周期范畴」章中作为函子修正工具。置闰把“不可公度性”这一周期分析中最顽固的数论困难，从被动接受的缺陷转化为主动利用的结构——这是本书对周期分析方法论的一项独立贡献。

\section*{本书的立场与限度}

周期数学并不先验地宣称所有被考察的系统都“具有周期结构”。它在每个应用章节中首先用三个判据提问（判据的完整表述见\textbf{引言}）。对于满足条件的系统，本书给出严格的周期结构描述；对于部分满足条件的系统，本书精确测量其偏离闭合轨道的程度；对于不满足条件的系统，本书明确指出其“周期性”只是现象层面的隐喻。这一诊断功能，是周期数学区别于一般跨学科类比的关键所在。

周期数学也不试图成为万有理论。它不解释为何某个特定的守恒量取某个特定的值，不推导某个具体系统的演化方程，不预测经济的未来走向。它提供的是一组诊断工具——一套统一的语法，用于严格地描述、比较和分析不同系统中出现的周期结构。它的成功不在于证明一切都是周期，而在于能够精确地指出：这个系统是周期的，因为它在某个群的作用下轨道闭合；那个系统看似有周期，实则缺乏守恒律的保证；这种积分变换和那种积分变换，是同一群上调和分析在不同参数化下的等价表示。\textbf{这些判定的正确性，需要逐章由具体的数学命题支撑。}

\section*{本书的旅程}

本书的旅程，从最简单的平移周期出发，逐步构建旋转周期与标度周期的数学理论，经由群表示论、调和分析、周期流形，最终抵达周期范畴的抽象统一。在此框架建立之后，我们将带着这套工具，依次穿越天文学、历法学、地理学、物理学、生物学、心理学、经济学、政治学、计算机科学和积分变换——在每一个领域中，\textbf{尽可能}以定理或严格命题的方式，揭示周期的数学结构如何照亮我们对自然与社会的理解；同时，\textbf{对于尚不能给出严格命题的部分，本书将明确标注为猜想或类比}。

封面上那幅嵌套的周转圆图景，正是这一旅程的起点。三个同心圆上的点，是最外层的天球；圆心的红点延伸出的嵌套小圆，是行星的轨迹；而封面的两行公式，则是这一切的数学语言。当我们合上这本书时，希望读者看到的将不再仅仅是周期性的函数，而是一个由振幅、频率、相位与周期共同编织的、生机勃勃的宇宙——\textbf{当然，这一切的前提是，我们始终记得：数学描述的对象是模型，而非宇宙本身}。

\section*{读者对象与预备知识}

本书面向四类读者：希望把系统的周期结构当作代数与几何对象来处理的数学工作者；日常处理周期现象、希望获得统一词汇的理工科研究生与研究人员；在实际工作中面对不可公度采样比与频谱泄漏的工程师与数据分析者；以及关心现代数学与传统符号体系（如《易经》）之关系的、具有一般数学素养的读者。

数学各章所需的预备知识为：线性代数（矩阵、特征值、转置、正交群）；数学分析（极限、连续、可微、一致收敛）；群表示论（群、同态、轨道、群作用、表示）；以及初等复分析（复平面、指数函数、全纯函数）。需要拓扑与微分几何的两章——「流形基础与纤维丛」与「周期流形」——在章首各自声明其前置要求，缺少这部分准备的读者初读时可跳过这两章，而不影响其余部分的逻辑线索。应用各章除前十一章外不要求更多数学，且每章都会重述其所用的概念；「概率论与分布」与「时序分析」两章另需概率论与数理统计的入门课程基础。

\textbf{阅读路线}：以数学内核为主要兴趣的读者可读第一章至第二十四章，以及第四十六章；以某一应用为主要兴趣的读者可读第一章、第二部分中相关的一章，再读所关心的应用章；只需计算工具的读者可从「离散采样逼近、同余角与置闰」「调和分析」「周转圆激活函数与函数逼近」「建模与分析：从数据到周转圆」四章入手。

\section*{体例与标注约定}

本书数学各章遵循如下顺序：\textbf{定义} $\to$ \textbf{命题} $\to$ \textbf{引理} $\to$ \textbf{定理} $\to$ \textbf{证明} $\to$ \textbf{推论} $\to$ \textbf{性质} $\to$ \textbf{例子} $\to$ \textbf{注释}。交叉引用一律按名称给出，全书的引用机制使用 \LaTeX{} 的 \texttt{cleveref} 宏包所提供的标准功能。

\textbf{每一条断言的证明状态都被标注}。凡在本书中给出证明者，以定义、命题、引理、定理、推论或性质的形式呈现，并附证明。凡在本书中未给出证明者，以下列三种标签之一标出：\textbf{猜想}，指陈述精确但证明尚未获得；\textbf{类比}，指两个结构之间的对应具有启发性，但不是双射；\textbf{研究纲领}，指一个用于组织一类问题、而非断言某条定理的框架。\textbf{这三种标签应被读作警示，而非修饰}：凡带有其中任一标签的论断，均无证明可用，且本书不在任何推导中把它当作前提。

记号：$\mathbb{N}, \mathbb{Z}, \mathbb{Q}, \mathbb{R}, \mathbb{C}, \mathbb{H}$ 分别表示自然数、整数、有理数、实数、复数与四元数。三种对称中心记作 $c=0$（加法反射）、$c=1$（乘法缩放反射）与 $c=i$（乘法取逆）。周期记作 $T$，对应的角频率记作 $\omega = 2\pi/T$，相位算子记作 $\mathcal{O}_k^{(T)} = e^{i2\pi k/T}$。周期卷积记作 $\circledast$；当需要显示阶数 $d$ 或窗核时，记作 $\mathcal{K}^{(d)}\{f,g\}$。

\textbf{关于史料}：序言与「易经系统之周期结构」章中的历史与文化材料分两类。关于传统历法推演、《易经》以及天文学史的部分，均以\textbf{传统说法或据称史料}引用；它们不是本书任何数学推导的前提，其具体数值应视为象征性设定。关于现代数学的部分，则连同其前提一并给出。凡需要原始文献而本书无法提供者，正文均标注为「史料待考」。

\section*{致谢}

本书起源于一个尝试：用同一套记号，把出现在天体力学、晶体学、信号处理与历法推演等相距甚远的领域中的周期结构写下来。在此过程中，它欠下了许多债。它首先欠那些定理被本书引用的数学家，以及参考文献中所列教材与专著的作者——本书中出现的群表示论、调和分析、数论与微分几何的标准结果都归他们所有，而本书中若有任何新意，也都是站在这些结果之上得到的。

它尤其欠中国古典数学与天文学文献——包括《淮南子》《皇极经世》、二十四节气与六十甲子——一份债：这些文献提供了本书所称“置闰”的公倍推理的最早实例。\textbf{需要再次强调}：本书给出的现代数学读解属于作者本人的诠释，不是对这些文本的历史断言。

感谢开源工具的维护者，没有它们本书无法完成：排版所用的 \TeX{} Live 与 \LaTeX{} 宏包，以及中文字体所采用的免费字形的开发者。仍然存在的错误，均属作者本人。

\textbf{尤其}欢迎更正、反例与反驳。由于本书是以研究纲领而非完成建筑的形式提出的，对它最有用的回应，是指出它在何处失败——并尽可能精确。

\vspace{1em}
\begin{flushright}
周期数学独立研究\\
2026 年 9 月
\end{flushright}

